% Generated draft: requires mathematical and source review.
\documentclass[11pt]{article}
\usepackage[margin=25mm]{geometry}
\usepackage{fontspec}
\setmainfont{Latin Modern Roman}
\usepackage{amsmath,amssymb,mathtools,unicode-math}
\setmathfont{Latin Modern Math}
\usepackage{enumitem,needspace,xurl,hyperref}
\setlength{\parindent}{0pt}
\setlength{\parskip}{5pt}
\setlength{\emergencystretch}{3em}
\newcommand{\N}{\mathbb{N}}
\newcommand{\Z}{\mathbb{Z}}
\newcommand{\Q}{\mathbb{Q}}
\newcommand{\R}{\mathbb{R}}
\newcommand{\C}{\mathbb{C}}
\newcommand{\F}{\mathbb{F}}
\newcommand{\A}{\mathbb{A}}
\newcommand{\PP}{\mathbb{P}}
\newcommand{\E}{\mathbb{E}}
\newcommand{\eps}{\varepsilon}
\newcommand{\dd}{\,\mathrm{d}}
\begin{document}
\section*{UnsolvedMath ID: 11000131}
\textbf{Source problem code:} catalogue entry
\subsection{Definitions and mathematical statement}
Let \(F\) be the source's surface, with its spaces \(\mathrm{PL}_0(F)\), \(\mathrm{PL}(F)\) and groups \(\mathrm{MC}(F)\), \(\mathrm{PMC}(F)\). Their definitions and surface conventions are not included. The group names conventionally denote mapping class and pure mapping class groups, but the exact boundary/puncture conditions require confirmation; the letters PL are not a definition of a particular lamination or piecewise-linear space.

Problem 1 asks to understand the non-Hausdorff quotients obtained from the alternatives
\[
 X/G,\qquad
 X\in\{\mathrm{PL}_0(F),\mathrm{PL}(F)\},\quad
 G\in\{\mathrm{MC}(F),\mathrm{PMC}(F)\},
\]
either classically or as quantum geometric objects. For an orbit-space interpretation the quotient topology declares a subset open exactly when its preimage under \(X\to X/G\) is open. The source does not supply the topology of \(X\), the precise actions, or whether the intended quantum object is attached to that orbit space or to richer action data. No quantization framework or criterion for ``understand'' is specified.

The context reports quantizations of \(T(F)\) in [5] and [17], surveyed respectively in [6] and [35], and a quantization of \(\mathrm{PL}_0(F^1_{1,0})\) in [6]. The surface-index convention in this flattened notation is missing. In contrast it cites [21,54] for a program treating real quadratic number fields through quantum-torus limits of elliptic curves, emphasizing quantization of limiting curves rather than Teichmüller space. These reports do not define a common quantum model for the requested non-Hausdorff quotients or assert that one of the existing constructions settles the problem.
\subsection{Short English statement}
Understand the specified non-Hausdorff mapping-class quotients, through classical topology or an appropriate quantum-geometric model. The spaces, actions and standard of understanding are missing. Existing Teichmüller and special-case quantizations, and a contrasting quantum-torus program, are retained as source context.
\subsection{Sources}
\begin{itemize}
\item Local UnsolvedMath registry, permanent ID \(11000131\), exact source code \texttt{catalogue entry}, Problem 1, PDF page 110. No individual proposer or proposed year is supplied.
\item Farb, edited volume \emph{Problems on Mapping Class Groups and Related Topics}: \url{https://www.math.uchicago.edu/~farb/papers/mcgbook.pdf}. The volume has 329 explicitly typeset numbered items across chapters, with repeated local numbering distinguished by PDF page.
\item Internal references: [5] and [17] for quantizing \(T(F)\), surveys [6] and [35] respectively; [6] for the displayed special \(\mathrm{PL}_0\) case; [21,54] for the real-quadratic-field/quantum-torus limiting-curve program. Complete bibliography entries are absent.
\item Extraction pdf-font-structure, accessed 2026-07-29. Source evidence NEEDS\_REVIEW presents an open problem and calls for a current primary-literature audit. Supplied OPEN-TRIAGE summary mentions Penner's decorated/PL Teichmüller theory and calls full understanding of the quotient quantum objects open, without precise further citations.
\item Dataset attribution: Ulam AI / UnsolvedMath Contributors, CC BY 4.0. Source rights NEEDS\_REVIEW: publicly accessible author/publisher copy, redistribution terms requiring review.
\end{itemize}
\end{document}
